\medskip

On se place dans un repère \Oij{} orthogonal.

\begin{enumerate}
\item On considère la fonction $g$ définie pour tout réel $x$ par 
$g(x) = x^2 - 5x + 4$.\\
On note $\mathcal{P}$ la courbe représentative de la fonction $g$.

	\begin{enumerate}
		\item 
$x^2-5x+4$ est un polynôme du second degré.

$\Delta=(-5)^2-4\times 4 = 25-16=9 = 3^2$

Le polynôme admet donc deux racines: 

$x'=\dfrac{-b-\sqrt{\Delta}}{2a}=\dfrac{5-3}{2}=1$ et $x'' =\dfrac{-b+\sqrt{\Delta}}{2a}=\dfrac{5+3}{2}=4$
		
On en déduit que $g(x)=(x-1)(x-4)$. On détermine alors le signe de $g$.		
		
		\[\begin{tablvar}[4em]{3}
\hline
x & -\infty && 1 && 4 &&  +\infty\\
\hline
x-1 &  & -  & \barre[0]& + & \barre&  + & \\
\hline
x-4 &  & -  & \barre& - & \barre[0] &  + & \\
\hline
\text{signe de } g &  & + & \barre[0]& - & \barre[0] &  + & \\
\hline
\end{tablvar}\]

		\item On considère un entier naturel $n$ quelconque.

On note $\text{A}_n$ le point de la courbe $\mathcal{P}$ d'abscisse $n$.

On note $a_n$ le coefficient directeur de la droite $(\text{A}_n\text{A}_{n+1})$ donc
$a_n=\dfrac{y_{\text{A}_{n+1}} - y_{\text{A}_n}}{x_{\text{A}_{n+1}} - x_{\text{A}_n}}$.

Par définition, $x_{\text{A}_n}=n$, et comme le point $\text{A}_n$ est un point de la courbe $\mathcal{P}$, on a:
$y_{\text{A}_n} = g\left (x_{\text{A}_n}\right )=g(n)=n^2-5n+4$.

De même:
$\begin{aligned}[t]
y_{\text{A}_{n+1}}
& = g\left (x_{\text{A}_{n+1}}\right )=g(n+1) =(n+1)^2-5(n+1)+4\\
& = n^2+2n+1-5n-5+4 = n^2-3n
\end{aligned}$

Donc:  $y_{\text{A}_{n+1}} - y_{\text{A}_{n}} = \left (n^2-3n\right ) - \left (n^2-5n+4\right ) = n^2 -3n -n^2 +5n -4=2n-4$.

Or $x_{\text{A}_{n+1}} - x_{\text{A}_{n}} = n+1-n=1$;
donc $a_n=\dfrac{y_{\text{A}_{n+1}} - y_{\text{A}_n}}{x_{\text{A}_{n+1}} - x_{\text{A}_n}} = \dfrac{2n-4}{1}=2n-4$.



		\item 
$a_{n+1}-a_n= \left (2(n+1)-4\right ) - \left (2n-4\right ) = 2n+2-4-2n+4=2$

Donc la suite $(a_n)$ est arithmétique de raison 2.
	\end{enumerate}
	
\item Soit $f$ la fonction définie pour tout réel $x$ de [0,5~;~8] par
$f(x) = x - 5 + \dfrac 4 x.$

On note $\mathcal{C}$ la courbe représentative de la fonction $f$.
	\begin{enumerate}
		\item 
$f(x) = x - 5 + \dfrac 4 x = \dfrac{x^2}{x} - \dfrac{5x}{x} + \dfrac{4}{x}= \dfrac{x^2-5x+4}{x}=\dfrac{g(x)}{x}$		

		\item 
La position de la courbe $\mathcal{C}$ par rapport à l'axe des abscisses dépend du signe de $f(x)$: si $f(x)>0$, alors la courbe $\mathcal{C}$ est au-dessus de l'axe des abscisses, et  si $f(x) < 0$, alors la courbe $\mathcal{C}$ est en dessous de l'axe des abscisses.

$f(x)=\dfrac{g(x)}{x}$, et sur $[0,5\;; 8]$, $x>0$; donc $f(x)$ est du même signe que $g(x)$ sur  $[0,5\;; 8]$.

\medskip

\begin{list}{\textbullet}{D'après le résultat de la question \textbf{1.a.}:}
\item sur l'intervalle $[0,5\;; 1[$, $g(x)>0$ donc $f(x)>0$, donc  la courbe $\mathcal{C}$ est au-dessus de l'axe des abscisses;
\item sur l'intervalle $]1\;; 4[$, $g(x)<0$ donc $f(x)<0$, donc la courbe  $\mathcal{C}$ est en dessous de l'axe des abscisses;
\item sur l'intervalle $]4\;; 8]$, $g(x)>0$ donc $f(x)>0$, donc la courbe  $\mathcal{C}$ est au-dessus de l'axe des abscisses.
\end{list}

		\item On admet que la fonction $f$ est dérivable sur l'intervalle [0,5~;~8].


Sur  l'intervalle [0,5~;~8], $f'(x)=1-0 +4\left (-\dfrac{1}{x^2}\right ) = 1-\dfrac{4}{x^2} = \dfrac{x^2-4}{x^2}= \dfrac{(x-2)(x+2)}{x^2}$.

		\item On détermine le signe de $f'(x)$ sur l'intervalle [0,5~;~8].

\[\begin{tablvar}{5}
\hline
x & 0,5 && 2 && 8 \\
\hline
x-2 & & - & \barre[0] & + & \\
\hline
x+2 & & + & \barre & + & \\
\hline
(x-2)(x+2) & & - & \barre[0] & + & \\
\hline
x^2 & & + & \barre & + & \\
\hline
f'(x)= \dfrac{(x-2)(x+2)}{x^2} & & - & \barre[0] & + & \rule[-12pt]{0pt}{35pt} \\
\hline
\end{tablvar}\]

On peut calculer les valeurs intéressantes de $f(x)$ pour compléter le tableau de variations.

$f(0,5)=0,5-5+\dfrac{4}{0,5}=0,5-5+8=3,5$;
$f(2)=2-5+\dfrac{4}{2}=2-5+2=-1$, et
$f(8)=8-5+\dfrac{4}{8}=8-5+0,5=3,5$

\[\renewcommand{\fleche}{\ncline[linewidth=1.2pt,arrowsize=2pt 3,nodesep=2.5pt]{->}}
\begin{tablvar}{5}
\hline
x & 0,5 && 2 && 8 \\
\hline
f'(x)= \dfrac{(x-2)(x+2)}{x^2} & & - & \barre[0] & + & \rule[-12pt]{0pt}{35pt} \\
\hline
\variations{\mil{\text{variations de }f} & \haut{3,5} && \bas{-1} && \haut{3,5}} 
\hline
\end{tablvar}\]
		
		\item On réalise un schéma de l'allure de la courbe $\mathcal{C}$ sur lequel on fait apparaitre les résultats des questions \textbf{2.b} et \textbf{2.d}.
		
		\begin{center}
\psset{xunit=1cm,yunit=1cm,labelFontSize=\scriptstyle,arrowsize=2pt 3}
\begin{pspicture}(-1,-2)(9,4)
\psgrid[unit=1cm,subgriddiv=5,gridlabels=0,gridcolor=gray,subgridcolor=lightgray](-1,-2)(9,4)
\psaxes{->}(0,0)(-0.99,-1.99)(9,3.99)
\psplot[linecolor=blue,linewidth=1.2pt]{0.5}{1}{x 5 sub 4 x div add}
\psplot[linecolor=red,linewidth=1.2pt]{1}{4}{x 5 sub 4 x div add}
\psplot[linecolor=blue,linewidth=1.2pt]{4}{8}{x 5 sub 4 x div add}
\psline[linecolor=blue,linestyle=dashed,linewidth=1.5pt](2,0)(2,-1)(0,-1)
\psline[linecolor=blue,linestyle=dashed,linewidth=1.5pt](8,0)(8,3.5)(0,3.5)
\psline[linecolor=blue,linestyle=dashed,linewidth=1.5pt](0.5,0)(0.5,3.5)
\uput[l](0,3.5){\blue\scriptsize$3,5$} \uput[d](0.5,0){\blue\scriptsize$0,5$} 
\end{pspicture}
\end{center}

	\end{enumerate}
\end{enumerate}
